\id{IRSTI 50.13.13}{}

\begin{header}
\swa{A.~Tulegulov, A.~Jumagaliyeva, A.~Sadu, В.~Yemilov, M.~Nurmagambetov, J.~Zhantlesov, I.~Dyussebayev}{SOFTWARE FOR SELECTIVE LASER MELTING METHOD FOR ADDITIVE TECHNOLOGIES}

\tsp{1}A.~Tulegulov\alink{https://orcid.org/0000-0002-1195-6919},
\tsp{1}A.~Jumagaliyeva\alink{https://orcid.org/0000-0001-8632-5209}\envelope,
\tsp{1}A.~Sadu\alink{https://orcid.org/0009-0006-6417-1628},
\tsp{1}В.~Yemilov\alink{https://orcid.org/0009-0005-3116-9890},
\tsp{1}M.~Nurmagambetov\alink{https://orcid.org/0000-0001-6439-3809},
\tsp{1}J.~Zhantlesov\alink{https://orcid.org/0009-0004-9294-8292},
\tsp{2}I.~Dyussebayev\alink{https://orcid.org/0000-0001-8514-8300}
\end{header}

\begin{affil}
\tsp{1}K.Kulazhanov Kazakh University of Technology and Business, Astana, Kazakhstan,

\tsp{2}K. Satbayev Kazakh National Research Technical University, Almaty, Kazakhstan

\corrauthor{Corresponding- author: jumagalievaainur.m@gmail.com}
\end{affil}

The article is devoted to the study of additive manufacturing
technologies, which have become increas-ingly important in recent years
in modern mechanical engineering, aerospace engineering, medicine, and
other high-tech industries. Of particular interest is the software for
selective laser melting (SLM) technology. Modern algorithms allow for
high dimensional accuracy, the ability to produce parts with specified
properties, and a shorter production cycle compared to traditional
manufacturing methods.

In this study, the focus is on the automated control system for printing
parts and assemblies using selective laser melting. The obtained data
can be further used in the development of algorithms for automated
selection of powder compositions and SLM modes.

The purpose of this work is to automate the technology of selecting the
composition of the powder material for selective laser melting. To
achieve this goal, the work involves studying the SLM technology and the
features of obtaining powders from technogenic waste, developing an
algorithm for automating the selection of powder composition, and
creating a software module that provides automatic selection of the
optimal powder composition and printing parameters.

The scientific novelty and significance of the work lie in the creation
of a software and hardware complex for the rapid selection of powder
composition and parameters of the selective laser melting process. The
implementation of the proposed approach makes it possible to produce
practical parts in a short time without compromising the quality and
performance of the material, and it also forms the basis for further
development of automated additive manufacturing control systems.

{\bfseries Keywords:} additive manufacturing, software, algorithm, software
and hardware complex, selective laser melting (SLM), automation.

\begin{header}
АДДИТИВТІ ТЕХНОЛОГИЯЛАРҒА АРНАЛҒАН СЕЛЕКТИВТІ ЛАЗЕРЛІК БАЛҚЫТУ ӘДІСІНІҢ БАҒДАРЛАМАЛЫҚ ЖАСАҚТАМАСЫ

\tsp{1}А.Д.~Тулегулов,
\tsp{1}А.М.~Джумагалиева\envelope,
\tsp{1}А.~Саду,
\tsp{1}Б.~Емилов,
\tsp{1}М.Ш.~Нурмагамбетов,
\tsp{1}Ж.~Жантлесов,
\tsp{2}И.М.~Дюсебаев
\end{header}

\begin{affil}
\tsp{1}К. Құлажанов атындағы Қазақ технология және бизнес университеті, Астана, Қазақстан,

\tsp{2}Қ.И. Сәтбаев атындағы Қазақ ұлттық техникалық зерттеу университеті, Алматы, Қазақстан,

e-mail: jumagalievaainur.m@gmail.com
\end{affil}

Мақала соңғы жылдары заманауи машина жасауда, аэроғарыштық техникада,
медицинада және басқа да жоғары технологиялық салаларда маңызды рөл
атқарған аддитивті өндіріс технологияларын зерттеуге арналған.
Селективті лазерлік балқыту (SLM) технологиясының бағдарламалық
жасақтамасы ерекше қызығушылық тудырады. Заманауи Алгоритмдер
өлшемдердің жоғары дәлдігін, берілген қасиеттері бар бөлшектерді жасау
мүмкіндігін және дәстүрлі өндіріс әдістерімен салыстырғанда өндіріс
циклін қысқартуға мүмкіндік береді.

Осы зерттеу шеңберінде талдау объектісі ретінде селективті лазерлік
балқыту әдісімен бөлшектер мен тораптарды басып шығару процесін
басқарудың автоматтандырылған жүйесі қарастырылады. Алынған мәліметтер
болашақта ұнтақтардың құрамын және SLM режимдерін автоматтандырылған
таңдау алгоритмдерін жасау кезінде қолданыла алады.

Бұл жұмыстың мақсаты селективті лазерлік балқыту үшін ұнтақ материалының
құрамын таңдау технологиясын автоматтандыру болып табылады. Жұмыста
қойылған мақсатқа жету үшін SLM технологиясын және техногендік
қалдықтардан ұнтақтарды алу ерекшеліктерін зерделеу, ұнтақ құрамын
таңдауды автоматтандыру алгоритмін әзірлеу, сондай-ақ ұнтақтың оңтайлы
құрамы мен басып шығару параметрлерін автоматты түрде таңдауды
қамтамасыз ететін бағдарламалық модуль құру көзделеді.

Жұмыстың ғылыми жаңалығы мен маңыздылығы ұнтақтың құрамын және
селективті лазерлік балқу процесінің параметрлерін жедел таңдау үшін
бағдарламалық-техникалық кешен құру болып табылады. Ұсынылған тәсілді
іске асыру материалдың сапасы мен пайдалану қасиеттерін жоғалтпай қысқа
мерзімде практикалық бөлшектерді дайындау мүмкіндігін қамтамасыз етеді,
сондай-ақ қосымша өндірісті басқарудың автоматтандырылған жүйелерін одан
әрі дамыту үшін негіз қалыптастырады.

{\bfseries Түйін сөздер}: аддитивті өндіріс, бағдарламалық қамтамасыз ету,
алгоритм, бағдарламалық-тех\-никалық кешен, селективті лазерлік балқыту
(SLM), автоматтандыру.

\begin{header}
ПРОГРАММНОЕ ОБЕСПЕЧЕНИЕ МЕТОДА СЕЛЕКТИВНОГО ЛАЗЕРНОГО ПЛАВЛЕНИЯ ДЛЯ АДДИТИВНЫХ ТЕХНОЛОГИЙ

\tsp{1}А.Д.~Тулегулов,
\tsp{1}А.М.~Джумагалиева\envelope,
\tsp{1}А.~Саду,
\tsp{1}Б.~Емилов,
\tsp{1}М.Ш.~Нурмагамбетов,
\tsp{1}Ж.~Жантлесов,
\tsp{2}И.М.~Дюсебаев
\end{header}

\begin{affil}
\tsp{1}Казахский университет технологии и бизнеса им. Кулажанова, Астана, Казахстан,

\tsp{2}Казахский национальный исследовательский технический университет им. К. Сатпаева, Алматы, Казахстан,

e-mail: jumagalievaainur.m@gmail.com
\end{affil}

Статья посвящена исследованию аддитивных технологий производства,
которые в последние годы приобрели значительную важность в современном
машиностроении, аэрокосмической технике, медицине и других
высокотехнологичных отраслях промышленности. Особый интерес представляет
программное обеспечение технологии селективного лазерного плавление
(SLM). Современные алгоритмы позволяют обеспечить высокую точность
размеров, возможность изготовления деталей с заданными свойствами и
сокращение производственного цикла по сравнению с традиционными методами
производства.

В рамках данного исследования в качестве объекта анализа рассматривается
автоматизированная система управления процессом печати деталей и узлов
методом селективного лазерного плавления. Полученные данные могут быть
использованы в дальнейшем при разработке алгоритмов автоматизированного
подбора состава порошков и режимов SLM.

Целью данной работы является автоматизация технологии подбора состава
порошкового материала для селективного лазерного плавления. Для
достижения поставленной цели в работе предполагается изучение технологии
SLM и особенностей получения порошков из техногенных отходов, разработка
алгоритма автоматизации подбора состава порошка, а также создание
программного модуля, обеспечивающего автоматический выбор оптимального
состава порошка и параметров печати.

Научная новизна и значимость работы заключается в создании
программно-технического комплекса для оперативного подбора состава
порошка и параметров процесса селективного лазерного плавления.
Реализация предлагаемого подхода обеспечивает возможность изготовления
практических деталей в сжатые сроки без потери качества и
эксплуатационных свойств материала, а также формирует основу для
дальнейшего развития автоматизированных систем управления аддитивным
производством.

{\bfseries Ключевые слова:} аддитивное производство, программное
обеспечение, алгоритм, программно-технический комплекс, селективная
лазерная плавка (SLM), автоматизация.

\begin{multicols}{2}
{\bfseries Introduction.} Current trends in the development of additive
manufacturing technologies are closely associated with the
implementation of automation and intelligent control of technological
processes. The automation of powder composition selection and selective
laser melting parameters reduces process preparation time, minimizes the
influence of the human factor, and ensures high reproducibility of
manufacturing results.

In this study, a gear manufactured by selective laser melting is
analyzed as an example of a functional component operating under
mechanical loads. To evaluate the quality of the formed material,
optical microscopy is employed to investigate the microstructure, the
characteristics of layer-by-layer formation, and the presence of defects
associated with the laser melting process. The obtained data can be
further used to develop algorithms for the automated selection of powder
compositions and SLM processing parameters.

The aim of this work is to automate the technology for selecting the
composition of powder materials for selective laser melting. To achieve
this aim, the study addresses the following objectives: investigation of
SLM technology and the features of powder production from industrial
waste; development of an algorithm for automating the selection of
powder composition; and creation of a software module that enables the
automatic determination of optimal powder compositions and SLM
processing parameters.

The scientific novelty of this research lies in the fact that, in the
Republic of Kazakhstan, studies focused on the development of composite
powder materials for additive manufacturing technologies, particularly
for selective laser melting, have been extremely limited. Within the
framework of this work, powder materials derived from locally available
industrial waste are obtained and analyzed for the first time, followed
by an assessment of the quality of manufactured parts based on
microstructural analysis {[}1{]}.

The scientific significance of the study consists in the development of
a domestically grounded methodology that enables the rapid selection of
powder compositions and selective laser melting parameters under various
technological conditions. The implementation of the proposed approach
makes it possible to manufacture functional parts within a short
timeframe without compromising material quality or performance
characteristics, and also establishes a foundation for the further
development of automated control systems in additive manufacturing.

{\bfseries Materials and мethods.} The object of research in this work is a
metal gear made by selective laser melting (Selective Laser Melting,
SLM) from powder material. The choice of a gear as an object of research
is due to the fact that this part belongs to the class of functional
elements of machines and mechanisms operating under conditions of cyclic
mechanical loads and increased requirements for geometric accuracy and
surface quality {[}2{]}.

Gears produced using the SLM process are of practical interest for
mechanical engineering and precision instrument manufacturing, as
additive technologies enable the fabrication of components with complex
geometries without the use of conventional machining operations. At the
same time, the quality of such products is largely determined by the
characteristics of microstructure formation during the layer-by-layer
laser melting of the powder material (fig.1a).

In this study, the gear is considered as a model component that allows
the assessment of the influence of additive manufacturing process
parameters on the structure of the formed material. Microstructural
analysis of the product makes it possible to identify characteristic
features of SLM technology, including the layered structure, structural
heterogeneity, and the presence of potential defects associated with
laser processing conditions and the properties of the initial powder
material (fig.1b).

For the microscopic examination of the gear structure, a MAGUS Metal 630
metallographic microscope equipped with a trinocular head was used. The
microscope is designed for the investigation of opaque metallic samples
in reflected light mode {[}3{]}. The optical system provides an
effective magnification range of approximately 50× to 500×, enabling the
acquisition of both overview images of the surface and detailed images
of local microstructural features. (fig.2)
\end{multicols}

\begin{figs}[Fig.1 - Infographics of SLM technology(а). The resulting
sample(b)]
    \fig[0.45\textwidth]{i/image87}[a]
    \fig[0.45\textwidth]{i/image88}[b]
\end{figs}

\begin{figs}[Fig.2 - Magus 630 Microscope]
    \fig[0.65\textwidth][5.5cm]{i/image89}
    \fig[0.3\textwidth][5.5cm]{i/image90}
\end{figs}

\begin{multicols}{2}
Microscopic observations were carried out in reflected light at various
magnifications, which made it possible to obtain a general view of the
gear surface as well as high-resolution images of selected areas
characterizing the microstructural features of the material {[}4{]}.

The microscope specifications are presented in Table 1.

The surface of the sample under investigation was subjected to visual
analysis without complex metallographic preparation, which is consistent
with the objectives of the work and the conditions of practical quality
control of additive manufacturing products. This approach allows us to
evaluate the possibility of using optical microscopy as an element of
operational and potentially automated control of products obtained by
the SLM method. (tab.2)
\end{multicols}

\tcap{Table 1 - MAGUS Metal 630 microscope specifications}
\begin{longtblr}[
  label = none,
  entry = none,
]{
  width = 0.9\linewidth,
  colspec = {Q[273]Q[721]},
  cells = {c},
  cells = {font = \small},
  row{1} = {font=\bfseries\small},
  hlines,
  vlines,
}
Parameter             & Value                                                                        \\
Microscope type       & Metallographic optical                                                       \\
Purpose               & Reflected light microscopy for metals and alloys                             \\
Image capture method  & Digital camera (MAGUS CBF30)                                                 \\
Magnification (total) & 50–500× (basic configuration)                                                \\
Eyepieces             & 10×/22 mm (10 mm interpupillary distance) (other eyepiece options available) \\
Objectives (included) & PL L5×/0,12; PL L10×/0,25; PL L20×/0,40; PL L50×/0,70                        
\end{longtblr}

\tcap{Table 2 -Main materials and methods of additive manufacturing of metals}
\begin{longtblr}[
  label = none,
  entry = none,
]{
  width = \linewidth,
  colspec = {Q[102]Q[115]Q[132]Q[102]Q[179]Q[306]},
  cells = {c},
  cells = {font = \small},
  row{1} = {font=\bfseries\small},
  cell{2}{1} = {r=4}{},
  hlines,
  vlines,
}
Category         & Material / Method   & Ultimate tensile strength, MPa & Yield strength, MPa & Application area                          & Key advantages                                                                \\
Aluminum alloys  & AlSi10Mg            & 460 ± 20                       & 270 ± 10            & Aerospace, automotive, defense industries & High strength, corrosion resistance, relatively low cost                      \\
                 & AlSi12              & 409 ± 20                       & 211 ± 20            & Lattice and lightweight structures        & High bending strength, good castability                                       \\
                 & AlSi9Cu3            & 415 ± 15                       & 236 ± 8             & Aerospace and automotive components       & Improved mechanical properties compared to casting                            \\
                 & AlCuMgZr            & 451 ± 4                        & 446 ± 4             & High-load aerospace components            & Near-full density, high fatigue strength                                      \\
Magnesium alloys & WE43 (Mg–Nd–Zn–Zr)  & -                              & -                   & Biomedical implants                       & Biocompatibility, biodegradability, porous structure supporting bone ingrowth \\
Titanium alloys  & Ti-6Al-4V / Ti-45Nb & -                              & -                   & Implants, aerospace industry              & High specific strength, reduced Young’s modulus (Ti-45Nb)                     
\end{longtblr}

\begin{multicols}{2}
During the study, microphotographs of the gear surface were taken,
reflecting the nature of the layer-by-layer formation of the material,
the presence of porosity, inhomogeneities, and other structural features
{[}5-6{]}. The images obtained were used for subsequent analysis and
interpretation of the microstructure of the product, as well as for
comparing the observed features with data presented in scientific
publications on selective laser melting (fig.3).
\end{multicols}

\fig{i/image91}[Fig.3 - LPBF Operational Principles and Schematic.]

\begin{multicols}{2}
Figure 4 illustrates the schematic representation of the selective laser
melting (SLM) process used for the fabrication of metal components from
powder materials. The process is based on the layer-by-layer melting of
a thin powder layer by a high-energy laser beam, followed by
solidification and bonding with the previously formed layer.

During the manufacturing process, a powder dispenser supplies a
controlled amount of metal powder onto the build platform. A recoater
blade spreads the powder into a uniform layer with a predefined
thickness. The build chamber is filled with an inert argon atmosphere to
prevent oxidation of the molten material and to ensure stable melting
conditions {[}7-9{]}.

A fiber laser source generates a focused laser beam, which is directed
onto the powder bed by a scanning galvanometer system. The laser
selectively melts the powder according to the predefined scan strategy,
forming a localized molten pool. After laser exposure, the molten
material rapidly solidifies, resulting in the formation of a dense solid
layer.

{\bfseries Results and Discussions.} As a result of microscopic analysis of
the surface of a gear manufactured using selective laser melting,
micrographs were obtained that reflect the peculiarities of material
structure formation during the additive manufacturing process. The
investigation was conducted at various magnifications, which made it
possible to evaluate both the overall surface morphology of the product
and local structural features (fig.4).

\fig[0.4\textwidth]{i/image92}[Fig.4 - Photograph of a gear with 40x magnification]

At a magnification of 40× (fig.5), a heterogeneous surface structure
formed as a result of the layer-by-layer laser melting of the powder
material can be clearly observed {[}10{]}. The surface morphology is
characterized by alternating regions with different degrees of fusion,
which is a typical feature of components produced using the SLM process.
\end{multicols}

\fig[0.55\textwidth]{i/image93}[Fig.5 - Powder Bed Fusion Physical Manufacturing Cycle]

\begin{multicols}{2}
The arrangement of structural features indicates a pronounced
directionality of the material formation process, determined by the
laser scanning strategy and the sequence of layer deposition. Similar
structural characteristics are widely reported in scientific studies
devoted to the selective laser melting of metallic powders and confirm
the validity of the selected technological approach. (fig.6a). At a
magnification of 50× (fig.6b), a more detailed analysis of the surface
structure becomes possible. The micrographs reveal characteristic
boundaries between individual layers formed during the successive
melting of the powder material. The presence of pronounced layering is a
consequence of the discrete nature of the additive manufacturing process
and is determined by the applied layer thickness and laser processing
parameters.
\end{multicols}

\begin{figs}[Fig.6 - Photo of the gear from the reverse side 50×]
    \fig[0.45\textwidth]{i/image94}[a]
    \fig[0.45\textwidth]{i/image95}[b]
\end{figs}

\begin{multicols}{2}
In certain areas of the surface, areas with increased roughness and
local structural in homogeneities were identified. These features may be
associated with the uneven distribution of heat flows in the laser
exposure zone, as well as with variations in the properties of the
powder material, including particle size and shape {[}11{]}.

It should be noted that no pronounced macrodefects, such as large cracks
or regions of incomplete fusion, were detected in the investigated
sample. This indicates sufficient stability of the selective laser
melting process and appropriately selected baseline printing parameters
for the fabrication of a functional component.

Individual micrographs (fig.7) show the presence of small pores
distributed unevenly across the surface of the product. Porosity is one
of the most common defects in additive manufacturing using the SLM
method and can result from insufficient laser energy density, the
presence of gas inclusions in the powder, or incomplete fusion of
particles.

The size and morphology of the observed pores indicate microporosity,
which generally has a less pronounced effect on the operational
properties of the component compared to larger defects. Nevertheless,
the presence of even fine porosity should be taken into account in the
further selection of powder material composition and optimization of
selective laser melting process parameters.

The obtained results confirm the feasibility of using microscopic
analysis as a tool for the primary quality assessment of components
manufactured by selective laser melting, particularly at the stages of
powder composition development and process parameter optimization.

\fig[0.8\columnwidth]{i/image96}[Fig.7 - Morphology of gear wear
particles revealed by optical microscopy]
\vspace{-1em}
The obtained data confirm the feasibility of using microscopic analysis
as a tool for the preliminary quality assessment of components
manufactured by selective laser melting, particularly at the stages of
powder composition development and optimization of technological process
parameters.

Figure 8 presents optical micrographs of the surface of the gear
manufactured by selective laser melting at different magnifications. At
a magnification of 40×, a heterogeneous surface structure associated
with the layer-by-layer nature of material formation is clearly
observed. At a magnification of 500×, local structural inhomogeneities
and individual pores characteristic of SLM-fabricated components are
revealed.
\end{multicols}

\fig[0.7\textwidth]{i/image97}[Fig.8 - Photograph of sample images of a gear]

\begin{multicols}{2}
{\bfseries Conclusion.} During the study, the microstructural features of a
gear manufactured by selective laser melting were evaluated using
optical microscopy. The obtained micrographs revealed characteristic
features of the additive manufacturing process, including layer-by-layer
material formation, microstructural heterogeneity, and the presence of
microporosity.

The microstructural analysis showed that the investigated sample
exhibits structural features typical of SLM-produced components, while
no significant macrodefects, such as cracks or regions of incomplete
fusion, were detected. This indicates sufficient stability of the
selective laser melting process and an acceptable level of material
quality for functional applications.

The results of the study confirm that optical microscopy is an effective
and accessible method for the preliminary quality assessment of parts
manufactured using selective laser melting. This method enables the
rapid identification of key microstructural features and potential
defects, which is particularly important during the stages of powder
material selection and optimization of SLM processing parameters.

The obtained results can be used as a basis for the further development
of automated algorithms for selecting powder compositions and selective
laser melting parameters, as well as for the creation of quality control
systems in additive manufacturing aimed at improving the reproducibility
and reliability of manufactured components.

\emph{{\bfseries Funding.} This study is funded by the Committee on science
and higher education (No. AR23490424 Development of an additive
installation for the manufacture of metal objects for the defense
industry").}
\end{multicols}

\begin{center}
{\bfseries References}
\end{center}

\begin{refs}
1. Liu J., et al. Automated, economical, and environmentally-friendly
asphalt mix design based on machine learning and multi-objective grey
wolf optimization// Journal of traffic and transportation engineering.
-2024.
-\href{file:///C:/Users/admin/OneDrive/Desktop/ZU/2026/ARTICLE/Vol.\%2011(3)}{Vol.
11(3).} - P.381- 405.
\href{https://doi.org/10.1016/j.jtte.2023.10.002}{DOI
10.1016/j.jtte.2023.10.002}.

2. Sharov K, et al. The dry mortar consumption calculation automation in
the finishing work organization// IOP Conference Series: Materials
Science and Engineering. -2020. -Vol.~913:032042. DOI
\href{https://iopscience.iop.org/article/10.1088/1757-899X/913/3/032042}{10.1088/1757-899X/913/3/032042}.

3. Barbhuiya S., Saeed Sh. Artificial Intelligence in Concrete Mix
Design: Advances, Applications and Challenges// International Conference
on Innovation and Intelligence for Informatics, Computing, and
Technologies. -2023. -P.32-37. DOI
\href{https://doi.org/10.1109/3ICT60104.2023.10391485}{10.1109/3ICT60104.2023.10391485}.

4. Smakova N., Tulegulov A., et al. Software development for SLM-3D
printer prototype control system using klipper firmware program// KazUTB
Bulletin. -2025. -T.3(28). -P.166-176.
\href{https://doi.org/10.58805/kazutb.v.3.28-1138}{DOI
10.58805/kazutb.v.3.28-1138}.

5. Shegetaeva A., Smakova N., Tulegulov A., Steremhartz A. Аlgorithm for
multi-factored forecasting of network vulnerabilities: from CVE data
analysis to the development of a forecasting model// KazUTB Bulletin.
-2025. -T.2(27). -P.112-123. DOI
\href{https://doi.org/10.58805/kazutb.v.2.27-1006}{10.58805/kazutb.v.2.27-1006}.

6. Tulegulov A., Akishev K., et al. Evaluation of the use of automated
control systems for solving traffic control problems// KazUTB Bulletin.
-2024. -T.4(25). -P.8-15.
\href{https://doi.org/10.58805/kazutb.v.4.25-668}{DOI
10.58805/kazutb.v.4.25-668}.

7. Tulegulov A., Zhamangarin D., et al. The use of digital technologies
for intelligent traffic management//KazUTB Bulletin. -2024.
-T.2(23)-2024. -P.146-155. DOI
\href{https://doi.org/10.58805/kazutb.v.2.23-485}{10.58805/kazutb.v.2.23-485}.

8. Aryngazin K., Akishev K., Tulegulov A., et al. Development of an
information and logical database model for accounting and movement of
wind turbine equipment// KazUTB Bulletin. -2024. -T.1(22). -P.22-29. DOI
\href{https://doi.org/10.58805/kazutb.v.1.22-248}{10.58805/kazutb.v.1.22-248}.

9. Jumagaliyeva A., et al. The impact of blockchain and artificial
intelligence technologies in network security for e-voting//
International Journal of Electrical and Computer Engineering (IJECE).
-2024. -Vol.14 (6). -P.6723\textasciitilde6733. DOI 10.11591/ijece.
v14i6.pp.6723-6733.

10. Yergesh M., Tulegulov A., et al. Development of an intelligent
system automating managerial decision-making using big data//
Eastern-European Journal of Enterprise Technologies. -2023. -T.6. -№
3(126). -P.27-35 DOI 10.15587/1729-4061.2023.289395.

11. Bayzharikova М., Tulegulov A., et al. Evalution of the efficiency of
the technological process for the production of building products with
fillers from metallurgical slag// Metalurgija. -2024. -Vol.63(2).
-P.267-270. ISSN 0543-5846.
\end{refs}

\begin{info}
\hspace{1em}\emph{{\bfseries Information about the authors}}

Tulegulov A.D. - Candidate of Physical and Mathematical Sciences,
Associate Professor, K.Kulazhanov Kazakh University of Technology and
Business, Astana, Kazakhstan, e-mail:
tad62@ya.ru;

Jumagalieva A.M. - Master' s degree, Associate Professor,
K.Kulazhanov Kazakh University of Technology and Business, Astana,
Kazakhstan, e-mail:
jumagalievaainur.m@gmail.com;

Sadu A. - Master' s student, K.Kulazhanov Kazakh
University of Technology and Business, Astana, Kazakhstan, e-mail:
alishersadu9@gmail.com;

Yemilov Birzhan -- Master' s student, K.Kulazhanov Kazakh
University of Technology and Business, Astana, Kazakhstan, e-mail:
chuckbot30@gmail.com;

Nurmagambetov M.Sh. - Candidate of Agricultural Sciences, Associate
Professor, K.Kulazhanov Kazakh University of Technology and Business,
Astana, Kazakhstan, e-mail:
Nurmag@mail.ru;

Zhangabyl Kh.Zh. - Candidate of Physico-mathematical Sciences, Associate
Professor, K.Kulazhanov Kazakh University of Technology and Business,
Astana, Kazakhstan, e-mail:
Zhan@mail.ru;

Dyussebayev I. - PhD, K.I. Satpayev Kazakh National Research University,
Almaty, Kazakhstan, e-mail:
Nomad.i.m.13@mail.ru.

\hspace{1em}\emph{{\bfseries Сведения об авторах}}

Тулегулов А. Д.-кандидат физико-математических наук, ассоциированный
профессор, Казахский университет технологии и бизнеса им. К. Кулажанова,
Астана, Казахстан, e-mail: tad62@ya.ru;

Жумагалиева А. М.-магистр, ассоциированный профессор, Казахский
университет технологии и бизнеса им. К. Кулажанова, Астана, Казахстан,
e-mail: jumagalievaainur.m@gmail.com;

Саду А.-магистрант, Казахский университет технологии и бизнеса им.К.
Кулажанова, Астана, Казахстан, e-mail: alishersadu9@gmail.com;

Емилов Б.-магистрант, Казахский университет технологии и бизнеса им. К.
Кулажанова, Астана, Казахстан, e-mail: chuckbot30@gmail.com;

Нурмагамбетов М. Ш.-кандидат с.-х. наук, ассоциированный профессор,
Казахский университет технологии и бизнеса им. К. Кулажанова, Астана,
Казахстан, e-mail: Nurmag@mail.ru;

Жантлесов Ж.Х.-кандидат физико-математических наук, ассоциированный
профессор, Казахский университет технологии и бизнеса им. К. Кулажанова,
Астана, Казахстан, e-mail: Zhan@mail.ru;

Дюсебаев И.-доктор PhD, Казахский национальный исследовательский
университет им. К. Сатпаева, Алматы, Казахстан, e-mail: Nomad13@mail.ru.
\end{info}
